\section*{Capítulo \ref{ch:g12:arith} --- Aritmética}

\begin{solution}{exo:g12:arith:1}
$17 \times 119 = 2023$, luego $2026 = 17 \times 119 + 3$: cociente $119$ y
resto $3$. Para $-2026$: $-2026 = 17\times(-120) + 14$ (en efecto,
$17 \times 120 = 2040$ y $2040 - 2026 = 14$): cociente $-120$ y resto $14$
(el resto tiene que estar en $\intco{0}{17}$, así que \emph{no} es $-3$).
\end{solution}

\begin{solution}{exo:g12:arith:2}
\omterm{def:g12:complex:modulus}{Módulo} $10$: $7^2 = 49 \equiv 9 \equiv -1$. Por tanto,
$7^{100} = \left(7^2\right)^{50} \equiv (-1)^{50} = 1 \pmod{10}$: la última
cifra de $7^{100}$ es $1$.
\end{solution}

\begin{solution}{exo:g12:arith:3}
Euclides: $1071 = 2\times462 + 147$; $462 = 3\times147 + 21$;
$147 = 7\times21 + 0$. Luego $\gcd = 21$.

Sustitución hacia atrás:
$21 = 462 - 3\times147 = 462 - 3(1071 - 2\times462) =
7\times462 - 3\times1071$. Así que $u = -3$ y $v = 7$:
$1071\times(-3) + 462\times7 = 21$.
\end{solution}

\begin{solution}{exo:g12:arith:4}
Todo entero es $\equiv 0, 1, 2$ o $3 \pmod 4$ y, al elevar al cuadrado:
$0^2 \equiv 0$, $1^2 \equiv 1$, $2^2 = 4 \equiv 0$, $3^2 = 9 \equiv 1$. Así
que $n^2 \equiv 0$ o $1 \pmod 4$. Una suma de dos cuadrados es entonces
congruente con $0 + 0$, $0 + 1$ o $1 + 1$, es decir, con $0$, $1$ o
$2 \pmod 4$; nunca con $3$.
\end{solution}

\begin{solution}{exo:g12:arith:5}
\omterm{def:g12:arith:divides}{Divisibilidad} entre $2$: de $n$ y $n + 1$, uno es par. \omterm{def:g12:arith:divides}{Divisibilidad} entre
$3$: si $n \equiv 0$, entonces $3 \mid n$; si $n \equiv 1 \pmod 3$,
entonces $2n + 1 \equiv 3 \equiv 0$; y si $n \equiv 2$, entonces
$n + 1 \equiv 0$. En todos los casos, $3$ \omterm{def:g12:arith:divides}{divide} al producto. Como
$\gcd(2,3) = 1$, el \cref{cor:g12:arith:coprimeprod} da
$6 \mid n(n+1)(2n+1)$. (Esto vuelve a demostrar, además, que
$\frac{n(n+1)(2n+1)}{6}$, la suma de cuadrados del \cref{exo:g12:seq:1}, es
un entero.)
\end{solution}

\begin{solution}{exo:g12:arith:6}
Buscamos el inverso de $5$ \omterm{def:g12:complex:modulus}{módulo} $11$: probando (o con Bézout),
$5 \times 9 = 45 = 44 + 1 \equiv 1 \pmod{11}$. Multiplicando la \omterm{def:g12:arith:congruence}{congruencia}
por $9$:
\[
x \equiv 9 \times 3 = 27 \equiv 5 \pmod{11}.
\]
Las soluciones son los enteros $x = 5 + 11k$, con $k \in \Z$.
(Comprobación: $5\times5 = 25 \equiv 3 \pmod{11}$.)
\end{solution}

\begin{solution}{exo:g12:arith:7}
$\gcd(17, 40) = 1$, así que hay soluciones. Euclides:
$40 = 2\times17 + 6$; $17 = 2\times6 + 5$; $6 = 5 + 1$. Sustituyendo hacia
atrás: $1 = 6 - 5 = 6 - (17 - 2\times6) = 3\times6 - 17
= 3(40 - 2\times17) - 17 = 3\times40 - 7\times17$. Por tanto,
$17\times(-7) - 40\times(-3) = 1$: la solución particular
$(x_0, y_0) = (-7, -3)$.

Solución general de $17x - 40y = 1$: restando la relación particular,
$17(x + 7) = 40(y + 3)$; y como $\gcd(17, 40) = 1$, el lema de Gauss da
$40 \mid x + 7$, luego $x = -7 + 40k$ y después $y = -3 + 17k$, con
$k \in \Z$ (y todas ellas se comprueban).

Para $17x - 40y = 6$, multiplicamos la solución particular por $6$:
$(x_1, y_1) = (-42, -18)$, y el mismo razonamiento da
\[
x = -42 + 40k, \qquad y = -18 + 17k, \qquad k \in \Z .
\]
(Por ejemplo, con $k = 2$: $x = 38$, $y = 16$; en efecto,
$17\times38 - 40\times16 = 646 - 640 = 6$.)
\end{solution}

\begin{solution}{exo:g12:arith:8}
Supongamos $\sqrt2 = \frac ab$ con $a, b \in \N^*$; entonces
$a^2 = 2b^2$. En la factorización en \omterm{def:g12:arith:prime}{primos} de un cuadrado, todos los
exponentes son pares; así que el exponente de $2$ en $a^2$ es par, mientras
que en $2b^2$ es \omterm{def:g11:func:parity}{impar} (uno más que un número par). Dos factorizaciones del
mismo entero con exponentes distintos de $2$ contradicen la unicidad del
\cref{thm:g12:arith:fta}. Por tanto, no existe tal fracción:
$\sqrt2 \notin \Q$.
\end{solution}

\begin{solution}{exo:g12:arith:9}
\emph{1.} De $k\binom pk = p \binom{p-1}{k-1}$ se sigue que $p$ \omterm{def:g12:arith:divides}{divide} a
$k\binom pk$. Para $1 \leq k \leq p-1$, $p \nmid k$ y, siendo $p$ \omterm{def:g12:arith:prime}{primo},
$\gcd(p, k) = 1$, así que el lema de Gauss da $p \mid \binom pk$.

\emph{2.} Inducción sobre $a$. Para $a = 0$: $0^p \equiv 0$. Supongamos
$a^p \equiv a \pmod p$. Por el teorema del binomio,
\[
(a+1)^p = \sum_{k=0}^{p} \binom pk a^k
\equiv a^p + 1 \pmod p,
\]
porque todos los términos intermedios se anulan \omterm{def:g12:complex:modulus}{módulo} $p$ por el punto 1.
Por la hipótesis de inducción, $(a+1)^p \equiv a + 1 \pmod p$. Esto
demuestra $a^p \equiv a$ para todo $a \in \N$, y el caso $a < 0$ se sigue
escribiendo $a \equiv a + kp$ con un representante positivo adecuado.
\end{solution}

\begin{solution}{exo:g12:arith:10}
$n \equiv 2 \pmod 3$ y $n \equiv 3 \pmod 5$: escribimos $n = 2 + 3s$;
entonces $2 + 3s \equiv 3 \pmod 5$, es decir, $3s \equiv 1 \pmod 5$. El
inverso de $3$ \omterm{def:g12:complex:modulus}{módulo} $5$ es $2$ ($3\times2 = 6 \equiv 1$), luego
$s \equiv 2 \pmod 5$, digamos $s = 2 + 5t$, y $n = 8 + 15t$: las dos
primeras condiciones significan $n \equiv 8 \pmod{15}$.

Añadiendo $n \equiv 2 \pmod 7$: $8 + 15t \equiv 2 \pmod 7$, y
$15 \equiv 1 \pmod 7$, así que $t \equiv -6 \equiv 1 \pmod 7$, digamos
$t = 1 + 7u$. Por tanto, $n = 23 + 105u$:
\[
n \equiv 23 \pmod{105}.
\]
(Comprobación: $23 = 3\times7 + 2 = 5\times4 + 3 = 7\times3 + 2$.)
\end{solution}

\begin{solution}{pb:g12:arith:1}
\textbf{1.} $2026 = 289 \times 7 + 3$: $2026 \equiv 3 \pmod 7$. Potencias
de $7$ \omterm{def:g12:complex:modulus}{módulo} $10$: $7, 9, 3, 1$, con un ciclo de longitud $4$; y
$100 \equiv 0 \pmod 4$: la última cifra de $7^{100}$ es $1$.

\textbf{2.} $5^2 = 25 \equiv -1 \pmod{13}$, luego
$5^{116} = \left(5^2\right)^{58} \equiv (-1)^{58} = 1$ y
$5^{117} \equiv 5 \pmod{13}$.

\textbf{3.} El inverso de $3$ \omterm{def:g12:complex:modulus}{módulo} $7$ es $5$ ($15 \equiv 1$):
$x \equiv 5 \times 5 = 25 \equiv 4 \pmod 7$.

\textbf{4.} $97 = 2 \times 35 + 27$; $35 = 27 + 8$; $27 = 3 \times 8 + 3$;
$8 = 2 \times 3 + 2$; $3 = 2 + 1$. Sustituyendo hacia atrás:
$1 = 97 \times 13 + 35 \times (-36)$. Luego
$35 \times (-36) \equiv 1 \pmod{97}$: el inverso de $35$ es
$-36 \equiv 61 \pmod{97}$.

\textbf{5.} $a$ es invertible \omterm{def:g12:complex:modulus}{módulo} $n$ exactamente cuando
$\gcd(a, n) = 1$: Bézout proporciona $au + nv = 1$, es decir,
$au \equiv 1$; y, recíprocamente, la existencia de un inverso obliga al
\omterm{def:g12:arith:gcd}{máximo común divisor} a dividir a $1$.

\textbf{6.} $0{\cdot}10 + 3{\cdot}9 + 0{\cdot}8 + 6{\cdot}7 +
4{\cdot}6 + 0{\cdot}5 + 6{\cdot}4 + 1{\cdot}3 + 5{\cdot}2 +
2{\cdot}1 = 132 = 12 \times 11 \equiv 0 \pmod{11}$: válido.

\textbf{7.} La suma cambia en $wd$ con $1 \leq w \leq 10$ y
$1 \leq \abs d \leq 9$: como $11$ es \omterm{def:g12:arith:prime}{primo} y no \omterm{def:g12:arith:divides}{divide} a ninguno de los dos
factores, no puede dividir al producto
(\cref{thm:g12:arith:gauss} y \cref{prop:g12:arith:primedivides}): la suma
modificada nunca vuelve a ser $\equiv 0$, así que todo error de una sola
cifra dispara la alarma.

\textbf{8.} Intercambiar dos cifras adyacentes $a$ y $b$ (con pesos
$w + 1$ y $w$) cambia la suma en
$(w+1)b + wa - (w+1)a - wb = b - a \not\equiv 0$ si $a \neq b$: detectado.
El secreto es la \emph{primalidad} del $11$: \omterm{def:g12:complex:modulus}{módulo} $10$, productos como
$5 \times 2$ se anulan sin que ninguno de los factores sea nulo, así que un
error de $\pm 2$ en una posición de peso $5$ (o una trasposición
desafortunada) podría colarse.

\textbf{9.} Suma ponderada de las doce cifras: $119$; el dígito de control
tiene que completarla hasta un múltiplo de $10$: $1$ (código completo,
$978\,2940199\,051$). El EAN se pierde las trasposiciones adyacentes con
$2(a - b) \equiv 0 \pmod{10}$, es decir, con $\abs{a - b} = 5$:
intercambiar un $2$ y un $7$, por ejemplo, pasa desapercibido; el precio
del amable \omterm{def:g12:complex:modulus}{módulo} $10$.

\textbf{10.} Doblando una cifra de cada dos desde la derecha y plegando
($16 \to 7$, etc.), la suma da $80 \equiv 0 \pmod{10}$: la tarjeta de
prueba es válida.

\textbf{11.} Con un \omterm{def:g12:complex:modulus}{módulo} \omterm{def:g12:arith:prime}{primo}, todos los pesos son invertibles, así que
se detectan \emph{todos} los errores de una cifra y \emph{todas} las
trasposiciones adyacentes: el lujo del ISBN; los esquemas \omterm{def:g12:complex:modulus}{módulo} $10$
conservan cifras amables para las personas y aceptan un pequeño punto
ciego.

\textbf{12.} $2^{10} = 1024 = 3 \times 341 + 1 \equiv 1 \pmod{341}$, luego
$2^{340} = \left(2^{10}\right)^{34} \equiv 1$. Y, sin embargo,
$341 = 11 \times 31$ es compuesto: pasa el test de Fermat en base $2$ sin
ser \omterm{def:g12:arith:prime}{primo}. Moraleja: la \omterm{def:g12:arith:congruence}{congruencia} de Fermat es necesaria, no suficiente;
los tests de primalidad necesitan herramientas más finas (y las consiguen,
en los volúmenes universitarios).

\textbf{13.} $3d \equiv 1 \pmod{20}$: $d = 7$ ($21 = 20 + 1$).

\textbf{14.} $c = 4^3 = 64 \equiv 31 \pmod{33}$.

\textbf{15.} $31 \equiv -2$: $(-2)^7 = -128$, y
$-128 + 4 \times 33 = 4$: el texto cifrado se descifra como $m = 4$. La
cerradura gira.

\textbf{16.} \omterm{def:g12:complex:modulus}{Módulo} $3$: si $3 \nmid m$, entonces $m^2 \equiv 1$ (Fermat),
luego $m^{21} = m \cdot \left(m^2\right)^{10} \equiv m$; y si $3 \mid m$,
los dos miembros son $\equiv 0$. \omterm{def:g12:complex:modulus}{Módulo} $11$: $m^{10} \equiv 1$ o bien
$11 \mid m$, y $m^{21} = m \cdot \left(m^{10}\right)^2 \equiv m$. Tanto
$3$ como $11$ dividen a $m^{21} - m$ y, al ser \omterm{def:g12:arith:gcd}{primos entre sí}, su producto
$33$ también (Gauss): $m^{21} \equiv m \pmod{33}$. El exponente
$ed = 21 = 1 + 20k$ se construyó para que desaparecieran los dos exponentes
de Fermat ($2$ y $10$, que dividen a $20$).

\textbf{17.} Multiplicar dos \omterm{def:g12:arith:prime}{primos} de $300$ cifras cuesta un microsegundo;
recuperarlos a partir de su producto derrota a todos los algoritmos
conocidos y a todos los ordenadores del mundo: la cerradura es una calle de
sentido único. (Nuestro $n = 33$ es esa calle a escala de juguete,
transitable en los dos sentidos.)

\textbf{18.} Probando restos (o construyendo con Bézout):
$n \equiv 8 \pmod{15}$: los recuentos $8, 23, 38, 53, \dots$ Unicidad
\omterm{def:g12:complex:modulus}{módulo} $15$: dos soluciones se diferencian en un múltiplo de $3$ y de $5$
y, por tanto, de $15$ ($3$ y $5$ son \omterm{def:g12:arith:gcd}{primos entre sí}; Gauss). El general
con $1000$ soldados anuncia el ``$8$'' con tres formaciones rápidas: el
viejo truco del recuento.

\textbf{19.} $10 \equiv 1 \pmod 9$ da $10^k \equiv 1$, luego
$\sum d_k 10^k \equiv \sum d_k$: un número y la suma de sus cifras son
congruentes \omterm{def:g12:complex:modulus}{módulo} $9$ (y \omterm{def:g12:complex:modulus}{módulo} $3$). Y $10 \equiv -1 \pmod{11}$ da
$\sum d_k 10^k \equiv \sum (-1)^k d_k$: la regla alternada. Dos reglas de
la infancia, a una línea cada una.

\textbf{20.} Las \omterm{def:g12:arith:congruence}{congruencias} convirtieron los restos en una aritmética
(Parte I); Bézout acuñó los inversos que resuelven las \omterm{def:g12:arith:congruence}{congruencias}
lineales y el $d$ de RSA (preguntas 4 y 13); el pequeño teorema de Fermat
abrió y cerró la cerradura (preguntas 15 y 16); y pegar \omterm{def:g12:complex:modulus}{módulos} \omterm{def:g12:arith:gcd}{primos
entre sí} contó soldados y remató la demostración (preguntas 16 y 18).
Veredicto sobre Hardy: el más puro de los teoremas que él conocía custodia
hoy cada compra; la pureza, con tiempo, es lo más aplicable que hay.
\end{solution}
